\section{Background and Motivation}
\label{sec:background}

\subsection{Multipathing and Weighted Allocation}

Application communication performance depends on both network utilization and the division of capacity among flows. Multiple transfers can collide on one path while leaving other paths underutilized. Even at full utilization, equal flow rates can delay application progress when transfers have different roles. In data analytics, some shuffle transfers determine when a subsequent stage starts. In distributed training, tensor completion times determine the overlap between communication and computation. These dependencies motivate per-flow allocation weights.

Multipath transmission and weighted allocation address these two requirements. Multipathing distributes traffic across parallel paths to reduce collisions and hotspots. Conventional ECMP assigns a flow to a fixed path; finer-grained transports support path switching or packet spraying. Weighted allocation assigns the resulting capacity to logical flows according to application policy, giving critical transfers larger shares while maintaining service for other traffic.

The two mechanisms operate at different granularities. A flow's packets can consume capacity on several paths, while its weight specifies one aggregate bandwidth share. A shared datacenter therefore needs an allocator that enforces flow-level policies across the service provided by multipath forwarding.

\subsection{Challenges of Weighted Allocation over Multipath}

A per-flow allocator must enforce accurate shares, adapt to contention changes, and keep control overhead low. Multipathing adds three challenges.

\paragraph{Challenge 1: Enforcing aggregate shares across bottlenecks.}
A multipath flow's throughput is the sum of its path rates, \(x_i=\sum_p x_{i,p}\). Each bottleneck serves only the packets traversing it, while the allocation objective applies to the entire flow. Consider two parallel 10-Gbps paths. Flow A uses the first path, flow C uses the second, and flow B uses both, with logical-flow weights \(w_A:w_B:w_C=1:2:1\). A feasible aggregate allocation is \((x_A,x_B,x_C)=(5,10,5)\)~Gbps, with B carrying 5~Gbps on each path. Applying the same weights independently at each link instead gives A and C approximately 3.33~Gbps each and B approximately 13.33~Gbps in total. Both paths are fully utilized, but the aggregate shares differ from the flow weights. Link utilization and flow-level allocation require distinct control.

\paragraph{Challenge 2: Tracking contention changes.}
A flow can participate in several contention groups simultaneously. In the preceding example, suppose a new flow D, with weight 2, starts transmitting over both paths. The desired aggregate rates of A, B, C, and D become approximately 3.33, 6.67, 3.33, and 6.67~Gbps, respectively. Although A and C use separate bottlenecks, both must adjust to D's arrival. An allocator must track these changes in aggregate service as overlapping contention groups evolve.

\paragraph{Challenge 3: Containing path-level control costs.}
Maintaining separate congestion and allocation state for each of a flow's \(P_i\) paths requires \(O(P_i)\) state and path-specific feedback. Under packet spraying, the selected path can change with every packet and may be unobservable at the sender. The allocation signal must capture aggregate resource use with per-flow state and processing costs independent of path count.

We assess these requirements through weight-normalized allocation error, convergence time, and per-flow state and processing overhead.

\subsection{The Cost of Enforcing Flow-Level Shares}

Allocation mechanisms place control at switches, at a centralized allocator, or at network-assisted end hosts. Multipath traffic changes the costs of each approach.

Switch-based weighted scheduling directly controls service at individual output links. Fine-grained enforcement requires flow classification, scheduling state, and queue management at line rate. A multipath flow can traverse several bottlenecks, extending these requirements to multiple switches. The resulting deployment depends on the available hardware queues, memory, and scheduling resources.

Centralized allocation computes target rates from traffic demands, routes, and capacity constraints. A change in a flow's demand or path usage can affect several contention groups, requiring updates to flow-resource relationships, recomputation of allocations, and distribution of new rates. Flow concurrency and traffic dynamics determine the computation and communication load; the update interval determines how quickly rates track current contention.

Network-assisted distributed allocation supplies switch-generated congestion feedback to end hosts. This reduces centralized computation but requires data-plane instrumentation. With packet spraying, the feedback for a logical flow reflects service at different paths and bottlenecks, so the sender must combine these observations into an aggregate allocation decision.

These approaches expose the value of a flow-level resource signal available from standard transport feedback. The next section establishes excess in-flight occupancy as such a signal and derives its relationship to bandwidth allocation.
